\documentclass[aspectratio=169,11pt]{beamer}

\usetheme{AnnArbor}
\usecolortheme{dolphin}
\setbeamertemplate{navigation symbols}{}
\setbeamertemplate{footline}[frame number]

\usepackage{fontspec}
\setmainfont{Tinos}
\setsansfont{Tinos}
\setmonofont[Scale=0.88]{DejaVu Sans Mono}
\usepackage{amsmath,amssymb,booktabs,array}
\usepackage{xcolor}
\usepackage{listings}
\usepackage{tikz}
\usepackage{ragged2e}

\definecolor{neublue}{RGB}{22,78,140}
\definecolor{neugreen}{RGB}{0,135,115}
\definecolor{lightblue}{RGB}{235,244,255}
\definecolor{lightgreen}{RGB}{235,250,245}
\definecolor{codebg}{RGB}{248,248,248}

\setbeamercolor{structure}{fg=neublue}
\setbeamercolor{title}{fg=neublue}
\setbeamercolor{frametitle}{fg=neublue}
\setbeamercolor{block title}{bg=neublue,fg=white}
\setbeamercolor{block body}{bg=lightblue,fg=black}
\setbeamercolor{block title alerted}{bg=red!70!black,fg=white}
\setbeamercolor{block body alerted}{bg=red!7,fg=black}
\setbeamercolor{block title example}{bg=neugreen,fg=white}
\setbeamercolor{block body example}{bg=lightgreen,fg=black}

\lstdefinestyle{py}{
  language=Python,
  backgroundcolor=\color{codebg},
  basicstyle=\ttfamily\scriptsize,
  keywordstyle=\color{neublue}\bfseries,
  stringstyle=\color{green!40!black},
  commentstyle=\color{gray!70!black}\itshape,
  showstringspaces=false,
  columns=fullflexible,
  keepspaces=true,
  frame=single,
  framerule=0.2pt,
  rulecolor=\color{gray!40},
  breaklines=true,
  tabsize=4
}

\lstdefinestyle{text}{
  backgroundcolor=\color{codebg},
  basicstyle=\ttfamily\scriptsize,
  showstringspaces=false,
  columns=fullflexible,
  keepspaces=true,
  frame=single,
  framerule=0.2pt,
  rulecolor=\color{gray!40},
  breaklines=true
}

\newcommand{\code}[1]{\texttt{#1}}
\newcommand{\key}[1]{\textbf{\textcolor{neublue}{#1}}}
\newcommand{\smallgap}{\vspace{0.4em}}

\title[Bài 3: Câu lệnh rẽ nhánh]{Bài 3: Câu lệnh rẽ nhánh}
\subtitle{Fundamental Programming Concepts in Python}
\author{Faculty of Data Science and Artificial Intelligence}
\institute{National Economics University}
\date{DSAI1003}

\begin{document}

% 1
\begin{frame}
  \titlepage
  \vspace{-1em}
  \begin{center}
    \small Dựa trên Chương 3 -- \textit{Decisions}, \textit{Python for Everyone}.
  \end{center}
\end{frame}

% 2
\begin{frame}{Vì sao cần câu lệnh rẽ nhánh?}
\begin{columns}[T]
\begin{column}{0.55\textwidth}
Ở Bài 2, chương trình chủ yếu chạy theo thứ tự từ trên xuống dưới.\smallgap

Trong thực tế, chương trình cần \key{đưa ra quyết định}: chọn hành động khác nhau tùy theo dữ liệu.
\end{column}
\begin{column}{0.40\textwidth}
\begin{exampleblock}{Ví dụ}
\begin{itemize}
  \item áp dụng giảm giá;
  \item kiểm tra dữ liệu nhập;
  \item phân loại điểm;
  \item chọn công thức tính thuế;
  \item xử lý chuỗi theo điều kiện.
\end{itemize}
\end{exampleblock}
\end{column}
\end{columns}
\end{frame}

% 3
\begin{frame}{Ý tưởng trung tâm}
\begin{center}
\Large
\key{Đánh giá điều kiện} $\rightarrow$ \key{Đúng/Sai} $\rightarrow$ \key{Chọn nhánh}
\end{center}
\vspace{1em}
\begin{block}{Luồng điều khiển}
\centering
\begin{tabular}{c c c}
Điều kiện đúng & $\Rightarrow$ & thực hiện hành động A \\
Điều kiện sai & $\Rightarrow$ & thực hiện hành động B
\end{tabular}
\end{block}
\vspace{1em}
Sau khi một nhánh được thực hiện, chương trình tiếp tục với các lệnh phía sau cấu trúc rẽ nhánh.
\end{frame}

% 4
\begin{frame}{Mục tiêu học tập}
Sau bài này, sinh viên có thể:
\begin{itemize}
  \item sử dụng \code{if}, \code{if/else}, \code{if/elif/else};
  \item dùng toán tử quan hệ để viết điều kiện;
  \item tổ chức khối lệnh bằng thụt lề đúng;
  \item thiết kế rẽ nhánh lồng nhau và nhiều phương án;
  \item sử dụng \code{and}, \code{or}, \code{not};
  \item thiết kế ca kiểm thử cho các nhánh và giá trị biên;
  \item phân tích chuỗi và kiểm tra dữ liệu đầu vào.
\end{itemize}
\end{frame}

% 5
\begin{frame}{Cấu trúc bài học}
\begin{columns}[T]
\begin{column}{0.48\textwidth}
\begin{enumerate}
  \item Câu lệnh \code{if}
  \item Toán tử quan hệ
  \item Rẽ nhánh lồng nhau
  \item Nhiều phương án
  \item Lưu đồ
\end{enumerate}
\end{column}
\begin{column}{0.48\textwidth}
\begin{enumerate}
\setcounter{enumi}{5}
  \item Ca kiểm thử
  \item Boolean và toán tử Boolean
  \item Phân tích chuỗi
  \item Kiểm tra dữ liệu đầu vào
  \item Tổng kết và bài tập
\end{enumerate}
\end{column}
\end{columns}
\end{frame}

\section{3.1 Câu lệnh if}

% 6
\begin{frame}[fragile]{Quyết định hai nhánh}
Ví dụ: tòa nhà không dùng số tầng 13. Nếu người dùng chọn tầng hiển thị lớn hơn 13, tầng thực tế giảm đi 1.
\smallgap
\begin{lstlisting}[style=py]
if floor > 13 :
    actualFloor = floor - 1
else :
    actualFloor = floor
\end{lstlisting}
\smallgap
\begin{block}{Điểm cần nhớ}
Chỉ \textbf{một nhánh} được thực hiện. Điều kiện đúng thì chạy nhánh \code{if}; điều kiện sai thì chạy nhánh \code{else}.
\end{block}
\end{frame}

% 7
\begin{frame}[fragile]{Cú pháp tổng quát}
\begin{columns}[T]
\begin{column}{0.50\textwidth}
\begin{block}{Có nhánh \code{else}}
\begin{lstlisting}[style=py]
if condition :
    statements1
else :
    statements2
\end{lstlisting}
\end{block}
\end{column}
\begin{column}{0.47\textwidth}
\begin{block}{Không có nhánh \code{else}}
\begin{lstlisting}[style=py]
if condition :
    statements
\end{lstlisting}
\end{block}
\end{column}
\end{columns}
\smallgap
Nhánh \code{else} là tùy chọn. Dùng khi cần xử lý riêng trường hợp điều kiện sai.
\end{frame}

% 8
\begin{frame}[fragile]{Câu lệnh ghép và khối lệnh}
\begin{lstlisting}[style=py]
if totalSales > 100.0 :
    discount = totalSales * 0.05
    totalSales = totalSales - discount
    print("You received a discount of", discount)
\end{lstlisting}
\begin{itemize}
  \item Dòng tiêu đề kết thúc bằng dấu hai chấm \code{:}.
  \item Các lệnh trong cùng một khối phải thụt lề cùng mức.
  \item Khối kết thúc khi mức thụt lề quay lại mức bên ngoài.
\end{itemize}
\begin{alertblock}{Quan trọng}
Trong Python, thụt lề là một phần của cú pháp, không chỉ để trình bày đẹp hơn.
\end{alertblock}
\end{frame}

% 9
\begin{frame}[fragile]{Ví dụ: mô phỏng thang máy}
\begin{lstlisting}[style=py]
##
# This program simulates an elevator panel that skips floor 13.
#

floor = int(input("Floor: "))

if floor > 13 :
    actualFloor = floor - 1
else :
    actualFloor = floor

print("The elevator will travel to the actual floor", actualFloor)
\end{lstlisting}
\begin{block}{Ví dụ chạy}
\begin{lstlisting}[style=text]
Floor: 20
The elevator will travel to the actual floor 19
\end{lstlisting}
\end{block}
\end{frame}

% 10
\begin{frame}[fragile]{Lỗi thường gặp: tab và thụt lề}
Python yêu cầu các lệnh trong cùng khối có mức thụt lề nhất quán.
\begin{lstlisting}[style=py]
if score >= 50 :
    print("Pass")
    print("Continue to the next course")
\end{lstlisting}
\begin{alertblock}{Tránh lỗi}
Không trộn tab và khoảng trắng. Nên cấu hình editor để phím Tab tạo 4 khoảng trắng.
\end{alertblock}
\end{frame}

% 11
\begin{frame}[fragile]{Mẹo: tránh lặp mã trong các nhánh}
\begin{columns}[T]
\begin{column}{0.48\textwidth}
\begin{alertblock}{Kém hơn}
\begin{lstlisting}[style=py]
if floor > 13 :
    actualFloor = floor - 1
    print("Actual floor:", actualFloor)
else :
    actualFloor = floor
    print("Actual floor:", actualFloor)
\end{lstlisting}
\end{alertblock}
\end{column}
\begin{column}{0.48\textwidth}
\begin{exampleblock}{Tốt hơn}
\begin{lstlisting}[style=py]
if floor > 13 :
    actualFloor = floor - 1
else :
    actualFloor = floor

print("Actual floor:", actualFloor)
\end{lstlisting}
\end{exampleblock}
\end{column}
\end{columns}
\end{frame}

% 12
\begin{frame}[fragile]{Chủ đề bổ sung: biểu thức điều kiện}
Python có dạng rút gọn cho một số quyết định đơn giản.
\begin{lstlisting}[style=py]
actualFloor = floor - 1 if floor > 13 else floor
\end{lstlisting}
\smallgap
\begin{block}{Khi nào dùng?}
\begin{itemize}
  \item Khi hai nhánh đều là biểu thức ngắn.
  \item Không nên dùng nếu làm mã khó đọc.
\end{itemize}
\end{block}
\end{frame}



% Added example
\begin{frame}[fragile]{Ví dụ bổ sung: miễn phí giao hàng}
Một cửa hàng miễn phí giao hàng nếu đơn hàng từ 500{,}000 đồng trở lên.
\begin{lstlisting}[style=py]
orderValue = float(input("Order value: "))

if orderValue >= 500000 :
    shippingFee = 0
else :
    shippingFee = 30000

total = orderValue + shippingFee
print("Shipping fee:", shippingFee)
print("Total:", total)
\end{lstlisting}
\begin{block}{Điểm học được}
Điều kiện \code{orderValue >= 500000} chia bài toán thành hai trường hợp rõ ràng: được miễn phí và không được miễn phí.
\end{block}
\end{frame}

\begin{frame}{Quiz 3.1: \texttt{if/else}}
Chọn đáp án đúng.
\begin{enumerate}
  \item Với \code{x = 10}, điều kiện \code{x > 10} có giá trị gì?
  \item Dấu \code{:} trong dòng \code{if condition :} dùng để làm gì?
  \item Nếu điều kiện của \code{if} sai và không có \code{else}, chương trình làm gì?
\end{enumerate}
\vspace{0.5em}
\begin{exampleblock}{Thảo luận nhanh}
Hãy dự đoán trước khi chạy code. Đây là thói quen quan trọng khi học lập trình.
\end{exampleblock}
\end{frame}

\begin{frame}{Đáp án Quiz 3.1}
\begin{enumerate}
  \item \code{False}, vì \code{10} không lớn hơn \code{10}.
  \item Dấu \code{:} báo hiệu bắt đầu một khối lệnh.
  \item Chương trình bỏ qua khối \code{if} và tiếp tục với lệnh phía sau.
\end{enumerate}
\end{frame}

\section{3.2 Toán tử quan hệ}

% 13
\begin{frame}{Các toán tử quan hệ}
\begin{center}
\begin{tabular}{ll}
\toprule
Toán tử & Ý nghĩa \\
\midrule
\code{<}  & nhỏ hơn \\
\code{<=} & nhỏ hơn hoặc bằng \\
\code{>}  & lớn hơn \\
\code{>=} & lớn hơn hoặc bằng \\
\code{==} & bằng nhau \\
\code{!=} & khác nhau \\
\bottomrule
\end{tabular}
\end{center}
\vspace{0.7em}
Kết quả của phép so sánh là một giá trị Boolean: \code{True} hoặc \code{False}.
\end{frame}

% 14
\begin{frame}[fragile]{Phép gán và phép so sánh bằng}
\begin{columns}[T]
\begin{column}{0.46\textwidth}
\begin{block}{Phép gán}
\begin{lstlisting}[style=py]
floor = 13
\end{lstlisting}
Lưu giá trị \code{13} vào biến \code{floor}.
\end{block}
\end{column}
\begin{column}{0.50\textwidth}
\begin{block}{Phép kiểm tra bằng nhau}
\begin{lstlisting}[style=py]
if floor == 13 :
    print("No such floor")
\end{lstlisting}
Kiểm tra \code{floor} có bằng \code{13} hay không.
\end{block}
\end{column}
\end{columns}
\begin{alertblock}{Lưu ý}
Trong điều kiện, thường dùng \code{==}, không dùng \code{=}.
\end{alertblock}
\end{frame}

% 15
\begin{frame}[fragile]{So sánh số thực: không nên kiểm tra bằng tuyệt đối}
Số thực trong máy tính có thể chứa sai số làm tròn rất nhỏ.
\begin{lstlisting}[style=py]
from math import sqrt

r = sqrt(2.0)
print(r * r)
\end{lstlisting}
\smallgap
Kết quả có thể là một giá trị rất gần 2, nhưng không đúng bằng 2.
\smallgap
\begin{exampleblock}{Cách kiểm tra gần bằng}
\begin{lstlisting}[style=py]
EPSILON = 1E-14
if abs(r * r - 2.0) < EPSILON :
    print("Close enough")
\end{lstlisting}
\end{exampleblock}
\end{frame}

% 16
\begin{frame}[fragile]{So sánh chuỗi}
Python có thể so sánh chuỗi theo thứ tự từ điển.
\begin{lstlisting}[style=py]
if name < "M" :
    print("First half of the directory")
else :
    print("Second half of the directory")
\end{lstlisting}
\begin{itemize}
  \item So sánh diễn ra từng ký tự từ trái sang phải.
  \item Chữ hoa và chữ thường có mã ký tự khác nhau.
  \item Nên chuẩn hóa chuỗi bằng \code{upper()} hoặc \code{lower()} khi cần.
\end{itemize}
\end{frame}

% 17
\begin{frame}{Cách làm: cài đặt một câu lệnh \texttt{if}}
\begin{enumerate}
  \item Xác định điều kiện rẽ nhánh.
  \item Mô tả điều xảy ra khi điều kiện đúng.
  \item Mô tả điều xảy ra khi điều kiện sai.
  \item Kiểm tra toán tử quan hệ tại giá trị biên.
  \item Loại bỏ phần xử lý lặp lại.
  \item Kiểm thử cả hai nhánh.
  \item Ghép thành chương trình Python.
\end{enumerate}
\end{frame}

% 18
\begin{frame}[fragile]{Ví dụ minh họa: lấy phần giữa của chuỗi}
Bài toán: nhập một chuỗi và lấy phần giữa.
\smallgap
Nếu chuỗi có độ dài lẻ, lấy 1 ký tự giữa. Nếu độ dài chẵn, lấy 2 ký tự giữa.
\begin{lstlisting}[style=py]
string = input("Enter a string: ")
length = len(string)

if length % 2 == 1 :
    position = length // 2
    result = string[position]
else :
    position = length // 2
    result = string[position - 1] + string[position]

print("Middle:", result)
\end{lstlisting}
\end{frame}



% Added example
\begin{frame}[fragile]{Ví dụ bổ sung: kiểm tra tuổi xem phim}
Một rạp phim yêu cầu người xem từ 18 tuổi trở lên đối với phim nhãn 18+.
\begin{lstlisting}[style=py]
age = int(input("Age: "))

if age >= 18 :
    print("Ticket allowed")
else :
    print("Not allowed for this movie")
\end{lstlisting}
\begin{block}{Giá trị biên}
Với điều kiện \code{age >= 18}, cần kiểm tra ít nhất các giá trị \code{17}, \code{18}, \code{19}.
\end{block}
\end{frame}

\begin{frame}{Quiz 3.2: toán tử quan hệ}
Dự đoán kết quả \code{True}/\code{False}.
\begin{center}
\begin{tabular}{ll}
\toprule
Biểu thức & Kết quả? \\
\midrule
\code{7 == 7.0} & \\
\code{7 != 7} & \\
\code{10 <= 10} & \\
\code{"An" < "Binh"} & \\
\code{3 + 4 > 10} & \\
\bottomrule
\end{tabular}
\end{center}
\end{frame}

\begin{frame}{Đáp án Quiz 3.2}
\begin{center}
\begin{tabular}{ll}
\toprule
Biểu thức & Kết quả \\
\midrule
\code{7 == 7.0} & \code{True} \\
\code{7 != 7} & \code{False} \\
\code{10 <= 10} & \code{True} \\
\code{"An" < "Binh"} & \code{True} \\
\code{3 + 4 > 10} & \code{False} \\
\bottomrule
\end{tabular}
\end{center}
\end{frame}

\section{3.3 Rẽ nhánh lồng nhau}

% 19
\begin{frame}{Khi nào cần rẽ nhánh lồng nhau?}
Một số bài toán cần quyết định theo nhiều tầng.\smallgap

Ví dụ: tính thuế phụ thuộc vào:
\begin{itemize}
  \item tình trạng hôn nhân;
  \item mức thu nhập;
  \item ngưỡng thuế tương ứng với từng nhóm.
\end{itemize}
\smallgap
Khi đó, nhánh bên ngoài chọn nhóm chính, nhánh bên trong xử lý điều kiện chi tiết hơn.
\end{frame}

% 20
\begin{frame}[fragile]{Ví dụ rẽ nhánh lồng nhau}
\begin{lstlisting}[style=py]
if maritalStatus == "s" :
    if income <= 32000 :
        tax = 0.10 * income
    else :
        tax = 3200 + 0.25 * (income - 32000)
else :
    if income <= 64000 :
        tax = 0.10 * income
    else :
        tax = 6400 + 0.25 * (income - 64000)
\end{lstlisting}
\begin{block}{Cách đọc}
Trước hết kiểm tra nhóm chính; sau đó, trong từng nhóm, kiểm tra ngưỡng thu nhập.
\end{block}
\end{frame}

% 21
\begin{frame}{Mẹo: mô phỏng bằng tay}
\begin{block}{Hand-tracing}
Mô phỏng chương trình trên giấy: ghi giá trị biến sau từng câu lệnh và đánh dấu nhánh được thực hiện.
\end{block}
\begin{exampleblock}{Khi nào rất hữu ích?}
\begin{itemize}
  \item rẽ nhánh lồng nhau;
  \item biểu thức Boolean phức tạp;
  \item lỗi kết quả không rõ nguyên nhân;
  \item kiểm thử giá trị biên.
\end{itemize}
\end{exampleblock}
\end{frame}



% Added example
\begin{frame}[fragile]{Ví dụ bổ sung: đăng nhập đơn giản}
Một chương trình kiểm tra tên người dùng trước, sau đó mới kiểm tra mật khẩu.
\begin{lstlisting}[style=py]
userName = input("Username: ")
password = input("Password: ")

if userName == "admin" :
    if password == "12345" :
        print("Login successful")
    else :
        print("Wrong password")
else :
    print("Unknown username")
\end{lstlisting}
\begin{block}{Cách đọc}
Nhánh trong chỉ được xét khi nhánh ngoài đã đúng.
\end{block}
\end{frame}

\begin{frame}[fragile]{Quiz 3.3: truy vết rẽ nhánh lồng nhau}
Với đoạn code sau, chương trình in gì nếu \code{x = 5}, \code{y = 2}?
\begin{lstlisting}[style=py]
if x > 0 :
    if y > 0 :
        print("A")
    else :
        print("B")
else :
    print("C")
\end{lstlisting}
\vspace{0.5em}
\begin{exampleblock}{Câu hỏi thêm}
Kết quả thay đổi thế nào nếu \code{x = 5}, \code{y = -2}?
\end{exampleblock}
\end{frame}

\begin{frame}{Đáp án Quiz 3.3}
\begin{itemize}
  \item Nếu \code{x = 5}, \code{y = 2}: in \code{"A"}.
  \item Nếu \code{x = 5}, \code{y = -2}: in \code{"B"}.
  \item Nếu \code{x <= 0}: nhánh trong không được xét, chương trình in \code{"C"}.
\end{itemize}
\end{frame}

\section{3.4 Nhiều phương án}

% 22
\begin{frame}[fragile]{Câu lệnh \texttt{if/elif/else}}
Khi có nhiều phương án loại trừ lẫn nhau, dùng \code{elif} để tránh nhiều \code{if} độc lập.
\begin{lstlisting}[style=py]
if condition1 :
    statements1
elif condition2 :
    statements2
elif condition3 :
    statements3
else :
    statements4
\end{lstlisting}
\begin{block}{Nguyên tắc}
Python kiểm tra từ trên xuống. Nhánh đầu tiên có điều kiện đúng sẽ được thực hiện.
\end{block}
\end{frame}

% 23
\begin{frame}[fragile]{Ví dụ: phân loại động đất}
\begin{lstlisting}[style=py]
if richter >= 8.0 :
    description = "Most structures fall"
elif richter >= 7.0 :
    description = "Many buildings destroyed"
elif richter >= 6.0 :
    description = "Many buildings damaged"
elif richter >= 4.5 :
    description = "Damage to poorly constructed buildings"
else :
    description = "No destruction of buildings"
\end{lstlisting}
\begin{block}{Điểm mấu chốt}
Các ngưỡng phải được đặt theo thứ tự phù hợp. Với ví dụ này, kiểm tra từ lớn xuống nhỏ.
\end{block}
\end{frame}

% 24
\begin{frame}[fragile]{Thứ tự điều kiện rất quan trọng}
\begin{columns}[T]
\begin{column}{0.48\textwidth}
\begin{alertblock}{Sai ý định}
\begin{lstlisting}[style=py]
if richter >= 4.5 :
    description = "Damage"
elif richter >= 6.0 :
    description = "Many damaged"
elif richter >= 8.0 :
    description = "Most fall"
\end{lstlisting}
\end{alertblock}
\end{column}
\begin{column}{0.48\textwidth}
\begin{exampleblock}{Vì sao sai?}
Nếu \code{richter = 8.5}, điều kiện đầu tiên \code{richter >= 4.5} đã đúng, nên các nhánh sau không bao giờ được xét.
\end{exampleblock}
\end{column}
\end{columns}
\end{frame}

% 25
\begin{frame}[fragile]{\texttt{elif} khác gì nhiều \texttt{if} độc lập?}
\begin{columns}[T]
\begin{column}{0.48\textwidth}
\begin{block}{Loại trừ lẫn nhau}
\begin{lstlisting}[style=py]
if score >= 90 :
    grade = "A"
elif score >= 80 :
    grade = "B"
elif score >= 70 :
    grade = "C"
\end{lstlisting}
\end{block}
\end{column}
\begin{column}{0.48\textwidth}
\begin{block}{Nhiều kiểm tra độc lập}
\begin{lstlisting}[style=py]
if hasLicense :
    print("License OK")
if hasInsurance :
    print("Insurance OK")
if hasRegistration :
    print("Registration OK")
\end{lstlisting}
\end{block}
\end{column}
\end{columns}
\end{frame}



% Added example
\begin{frame}[fragile]{Ví dụ bổ sung: tính phí vận chuyển theo vùng}
\begin{lstlisting}[style=py]
zone = input("Zone A/B/C: ").upper()

if zone == "A" :
    fee = 15000
elif zone == "B" :
    fee = 25000
elif zone == "C" :
    fee = 40000
else :
    fee = -1
    print("Invalid zone")

if fee >= 0 :
    print("Shipping fee:", fee)
\end{lstlisting}
\begin{block}{Điểm học được}
\code{if/elif/else} phù hợp khi các phương án loại trừ lẫn nhau.
\end{block}
\end{frame}

\begin{frame}{Quiz 3.4: chọn cấu trúc phù hợp}
Nên dùng \code{elif} hay nhiều \code{if} độc lập?
\begin{enumerate}
  \item Phân loại điểm A/B/C/D/F.
  \item Kiểm tra một đơn hàng có: mã giảm giá, freeship, quà tặng.
  \item Phân loại mức động đất theo ngưỡng Richter.
  \item Kiểm tra username: đủ dài, không có khoảng trắng, chỉ chữ/số.
\end{enumerate}
\end{frame}

\begin{frame}{Đáp án Quiz 3.4}
\begin{enumerate}
  \item Dùng \code{if/elif/else}: các mức điểm loại trừ nhau.
  \item Dùng nhiều \code{if}: các ưu đãi có thể cùng xảy ra.
  \item Dùng \code{if/elif/else}: mỗi mức Richter thuộc một nhóm.
  \item Thường dùng nhiều điều kiện kết hợp hoặc nhiều \code{if}, vì cần kiểm tra nhiều lỗi khác nhau.
\end{enumerate}
\end{frame}

\section{3.5 Lưu đồ và ca kiểm thử}

% 26
\begin{frame}{Lưu đồ: công cụ mô tả quyết định}
Lưu đồ giúp nhìn rõ luồng điều khiển trước khi viết mã.\smallgap

\begin{block}{Các thành phần thường dùng}
\begin{itemize}
  \item điểm bắt đầu/kết thúc;
  \item bước xử lý;
  \item điều kiện rẽ nhánh;
  \item mũi tên thể hiện luồng đi.
\end{itemize}
\end{block}
\smallgap
Trong bài học này, lưu đồ chủ yếu dùng để mô tả quyết định hai nhánh và nhiều phương án.
\end{frame}

% 27
\begin{frame}{Tránh luồng điều khiển kiểu ``spaghetti''}
Một chương trình tốt nên có luồng điều khiển dễ theo dõi.
\begin{alertblock}{Dấu hiệu cần tránh}
\begin{itemize}
  \item quá nhiều nhánh lồng nhau không cần thiết;
  \item điều kiện trùng lặp hoặc chồng lấn khó hiểu;
  \item cùng một xử lý bị copy ở nhiều nhánh;
  \item thứ tự kiểm tra không rõ lý do.
\end{itemize}
\end{alertblock}
\end{frame}

% 28
\begin{frame}{Ca kiểm thử: kiểm tra từng nhánh}
\begin{block}{Bao phủ nhánh}
Mỗi nhánh quan trọng của chương trình cần có ít nhất một ca kiểm thử.
\end{block}
\begin{exampleblock}{Ví dụ với điều kiện \code{floor > 13}}
\begin{itemize}
  \item \code{floor = 20}: nhánh đúng.
  \item \code{floor = 10}: nhánh sai.
  \item \code{floor = 13}: giá trị biên.
  \item \code{floor = 14}: ngay sau biên.
\end{itemize}
\end{exampleblock}
\end{frame}

% 29
\begin{frame}{Giá trị biên}
Lỗi thường xuất hiện gần các ngưỡng so sánh.\smallgap

Nếu điều kiện là \code{score >= 50}, hãy kiểm tra:
\begin{center}
\begin{tabular}{cl}
\toprule
Giá trị & Ý nghĩa \\
\midrule
\code{49} & ngay dưới ngưỡng \\
\code{50} & đúng tại ngưỡng \\
\code{51} & ngay trên ngưỡng \\
\bottomrule
\end{tabular}
\end{center}
\smallgap
Thiết kế ca kiểm thử trước khi lập trình giúp phát hiện lỗi logic sớm hơn.
\end{frame}



% Added exercise slide
\begin{frame}{Ví dụ bổ sung: bảng ca kiểm thử cho phân loại điểm}
\begin{center}
\begin{tabular}{cll}
\toprule
Điểm & Trường hợp & Kết quả mong đợi \\
\midrule
\code{-1} & ngoài miền dưới & \code{Invalid} \\
\code{0} & biên dưới & \code{F} \\
\code{59} & ngay dưới 60 & \code{F} \\
\code{60} & biên D & \code{D} \\
\code{89} & ngay dưới A & \code{B} \\
\code{90} & biên A & \code{A} \\
\code{101} & ngoài miền trên & \code{Invalid} \\
\bottomrule
\end{tabular}
\end{center}
\begin{block}{Nhận xét}
Bảng kiểm thử nên bao gồm cả nhánh hợp lệ, nhánh không hợp lệ và các giá trị sát ngưỡng.
\end{block}
\end{frame}

\begin{frame}{Hoạt động nhanh: thiết kế ca kiểm thử}
Bài toán: vé trẻ em nếu tuổi dưới 12; vé người lớn nếu từ 12 đến 59; vé người cao tuổi nếu từ 60 trở lên. Tuổi âm là không hợp lệ.
\begin{block}{Yêu cầu}
Hãy đề xuất ít nhất 6 ca kiểm thử, gồm:
\begin{itemize}
  \item một ca không hợp lệ;
  \item các ca ngay dưới, đúng tại, và ngay trên mỗi ngưỡng;
  \item một ca đại diện cho từng nhóm vé.
\end{itemize}
\end{block}
\end{frame}

\section{3.7 Boolean và toán tử Boolean}

% 30
\begin{frame}[fragile]{Biến Boolean}
Một biến Boolean lưu một trong hai giá trị: \code{True} hoặc \code{False}.
\begin{lstlisting}[style=py]
isValid = score >= 0 and score <= 100

if isValid :
    print("Valid score")
else :
    print("Invalid score")
\end{lstlisting}
\begin{block}{Lợi ích}
Biến Boolean giúp đặt tên cho một điều kiện, làm chương trình dễ đọc hơn.
\end{block}
\end{frame}

% 31
\begin{frame}{Toán tử \texttt{and}}
\code{and} đúng khi \textbf{cả hai} điều kiện đều đúng.
\begin{center}
\begin{tabular}{ccc}
\toprule
A & B & A \code{and} B \\
\midrule
False & False & False \\
False & True  & False \\
True  & False & False \\
True  & True  & True \\
\bottomrule
\end{tabular}
\end{center}
\begin{exampleblock}{Ví dụ}
\code{score >= 0 and score <= 100}
\end{exampleblock}
\end{frame}

% 32
\begin{frame}{Toán tử \texttt{or}}
\code{or} đúng khi \textbf{ít nhất một} điều kiện đúng.
\begin{center}
\begin{tabular}{ccc}
\toprule
A & B & A \code{or} B \\
\midrule
False & False & False \\
False & True  & True \\
True  & False & True \\
True  & True  & True \\
\bottomrule
\end{tabular}
\end{center}
\begin{exampleblock}{Ví dụ}
\code{ch == "Y" or ch == "y"}
\end{exampleblock}
\end{frame}

% 33
\begin{frame}[fragile]{Toán tử \texttt{not} và thứ tự ưu tiên}
\code{not} đảo giá trị Boolean.
\begin{lstlisting}[style=py]
if not isValid :
    print("Invalid input")
\end{lstlisting}
\smallgap
Thứ tự ưu tiên thường gặp:
\begin{center}
\begin{tabular}{cl}
\toprule
Ưu tiên & Toán tử \\
\midrule
cao & toán tử quan hệ như \code{<}, \code{==} \\
 & \code{not} \\
 & \code{and} \\
thấp & \code{or} \\
\bottomrule
\end{tabular}
\end{center}
Dùng ngoặc để làm rõ ý định.
\end{frame}

% 34
\begin{frame}[fragile]{Lỗi thường gặp: nhầm \texttt{and} với \texttt{or}}
Muốn kiểm tra điểm nằm ngoài khoảng 0 đến 100:
\begin{columns}[T]
\begin{column}{0.48\textwidth}
\begin{alertblock}{Sai}
\begin{lstlisting}[style=py]
if score < 0 and score > 100 :
    print("Invalid")
\end{lstlisting}
Không có số nào vừa nhỏ hơn 0 vừa lớn hơn 100.
\end{alertblock}
\end{column}
\begin{column}{0.48\textwidth}
\begin{exampleblock}{Đúng}
\begin{lstlisting}[style=py]
if score < 0 or score > 100 :
    print("Invalid")
\end{lstlisting}
Một trong hai điều kiện xảy ra là đủ.
\end{exampleblock}
\end{column}
\end{columns}
\end{frame}

% 35
\begin{frame}[fragile]{Đánh giá ngắn mạch}
Python đánh giá \code{and} và \code{or} từ trái sang phải và dừng khi đã biết kết quả.
\begin{lstlisting}[style=py]
if quantity > 0 and price / quantity < 10 :
    print("Good price")
\end{lstlisting}
\smallgap
Nếu \code{quantity > 0} sai, biểu thức thứ hai không được đánh giá. Điều này tránh lỗi chia cho 0.
\smallgap
\begin{block}{Ý nghĩa}
Thứ tự các điều kiện có thể ảnh hưởng đến cả tính đúng đắn và độ an toàn của chương trình.
\end{block}
\end{frame}

% 36
\begin{frame}{Định luật De Morgan}
De Morgan giúp biến đổi biểu thức phủ định.
\begin{center}
\begin{tabular}{ccl}
\toprule
Biểu thức & & Tương đương \\
\midrule
\code{not (A and B)} & $\Leftrightarrow$ & \code{not A or not B} \\
\code{not (A or B)}  & $\Leftrightarrow$ & \code{not A and not B} \\
\bottomrule
\end{tabular}
\end{center}
\begin{exampleblock}{Ví dụ}
\code{not (score >= 0 and score <= 100)} tương đương với \code{score < 0 or score > 100}.
\end{exampleblock}
\end{frame}



% Added quiz
\begin{frame}{Quiz 3.5: Boolean}
Dự đoán kết quả nếu \code{score = 75} và \code{absent = 2}.
\begin{enumerate}
  \item \code{score >= 50 and absent <= 3}
  \item \code{score >= 80 or absent <= 3}
  \item \code{not (score < 50)}
  \item \code{score < 0 or score > 100}
\end{enumerate}
\vspace{0.5em}
\begin{exampleblock}{Câu hỏi thêm}
Điều kiện nào phù hợp để kiểm tra \code{score} nằm trong đoạn 0--100?
\end{exampleblock}
\end{frame}

\begin{frame}{Đáp án Quiz 3.5}
Với \code{score = 75}, \code{absent = 2}:
\begin{enumerate}
  \item \code{True}
  \item \code{True}
  \item \code{True}
  \item \code{False}
\end{enumerate}
Điều kiện kiểm tra điểm hợp lệ:
\begin{center}
\code{score >= 0 and score <= 100}
\end{center}
\end{frame}

\begin{frame}[fragile]{Ví dụ bổ sung: điều kiện nhận học bổng}
Sinh viên được xét học bổng nếu GPA từ 3.2 trở lên và số tín chỉ đạt ít nhất 15.
\begin{lstlisting}[style=py]
gpa = float(input("GPA: "))
credits = int(input("Credits: "))

if gpa >= 3.2 and credits >= 15 :
    print("Scholarship candidate")
else :
    print("Not eligible")
\end{lstlisting}
\begin{block}{Điểm học được}
\code{and} dùng khi tất cả điều kiện con đều phải đúng.
\end{block}
\end{frame}

\section{3.8 Phân tích chuỗi}

% 37
\begin{frame}[fragile]{Toán tử thành viên \texttt{in} và \texttt{not in}}
\begin{lstlisting}[style=py]
if "@" in email :
    print("May be an email address")

if " " not in userName :
    print("No spaces")
\end{lstlisting}
\begin{block}{Ý nghĩa}
\code{in} kiểm tra một chuỗi con có xuất hiện trong chuỗi lớn hơn hay không.
\end{block}
\end{frame}

% 38
\begin{frame}{Một số phương thức chuỗi hữu ích}
\begin{center}
\begin{tabular}{ll}
\toprule
Phương thức & Công dụng \\
\midrule
\code{s.startswith(t)} & kiểm tra bắt đầu bằng \code{t} \\
\code{s.endswith(t)} & kiểm tra kết thúc bằng \code{t} \\
\code{s.find(t)} & vị trí xuất hiện đầu tiên của \code{t} \\
\code{s.count(t)} & số lần xuất hiện của \code{t} \\
\code{s.lower()} & chuyển về chữ thường \\
\code{s.upper()} & chuyển về chữ hoa \\
\bottomrule
\end{tabular}
\end{center}
\end{frame}

% 39
\begin{frame}[fragile]{Kiểm tra đặc điểm của chuỗi}
\begin{lstlisting}[style=py]
if text.isdigit() :
    number = int(text)

if name.isalpha() :
    print("Letters only")

if code.isalnum() :
    print("Letters and digits only")
\end{lstlisting}
\begin{itemize}
  \item \code{isdigit()}: toàn chữ số.
  \item \code{isalpha()}: toàn chữ cái.
  \item \code{isalnum()}: chữ cái hoặc chữ số.
\end{itemize}
\end{frame}

% 40
\begin{frame}[fragile]{Ví dụ: khảo sát chuỗi con}
\begin{lstlisting}[style=py]
string = input("Enter a string: ")
substring = input("Enter a substring: ")

position = string.find(substring)

if position == -1 :
    print("Not found")
else :
    print("First occurrence starts at", position)
    print("Number of occurrences:", string.count(substring))
\end{lstlisting}
\begin{block}{Kết nối với rẽ nhánh}
Phương thức chuỗi thường tạo ra điều kiện để chương trình quyết định bước xử lý tiếp theo.
\end{block}
\end{frame}



% Added example + quiz
\begin{frame}[fragile]{Ví dụ bổ sung: kiểm tra mã sinh viên}
Giả sử mã sinh viên hợp lệ nếu bắt đầu bằng \code{"SV"} và phần sau là chữ số.
\begin{lstlisting}[style=py]
studentID = input("Student ID: ").upper()

if studentID.startswith("SV") and studentID[2:].isdigit() :
    print("Valid student ID")
else :
    print("Invalid student ID")
\end{lstlisting}
\begin{alertblock}{Lưu ý}
Cần cẩn thận nếu chuỗi quá ngắn. Khi dạy sâu hơn, nên bổ sung kiểm tra \code{len(studentID) >= 3}.
\end{alertblock}
\end{frame}

\begin{frame}{Quiz 3.6: phân tích chuỗi}
Với \code{s = "Python 3"}, dự đoán kết quả:
\begin{enumerate}
  \item \code{"Py" in s}
  \item \code{s.startswith("py")}
  \item \code{s.upper().startswith("PY")}
  \item \code{s.isalpha()}
  \item \code{" " in s}
\end{enumerate}
\end{frame}

\begin{frame}{Đáp án Quiz 3.6}
Với \code{s = "Python 3"}:
\begin{enumerate}
  \item \code{True}
  \item \code{False}, vì phân biệt hoa/thường.
  \item \code{True}, vì đã chuẩn hóa bằng \code{upper()}.
  \item \code{False}, vì có khoảng trắng và chữ số.
  \item \code{True}.
\end{enumerate}
\end{frame}

\section{3.9 Kiểm tra dữ liệu đầu vào}

% 41
\begin{frame}{Vì sao cần kiểm tra dữ liệu đầu vào?}
Người dùng có thể nhập dữ liệu ngoài phạm vi mong đợi.\smallgap

Nếu chương trình xử lý ngay mà không kiểm tra, có thể xảy ra:
\begin{itemize}
  \item kết quả sai;
  \item lỗi chạy chương trình;
  \item thông báo khó hiểu;
  \item dữ liệu không nhất quán.
\end{itemize}
\begin{block}{Nguyên tắc}
Kiểm tra tính hợp lệ trước khi xử lý chính.
\end{block}
\end{frame}

% 42
\begin{frame}[fragile]{Kiểm tra miền giá trị}
\begin{lstlisting}[style=py]
score = int(input("Score: "))

if score < 0 or score > 100 :
    print("Invalid score")
else :
    if score >= 50 :
        print("Pass")
    else :
        print("Fail")
\end{lstlisting}
\begin{block}{Mẫu xử lý}
Trước hết loại bỏ dữ liệu không hợp lệ; sau đó xử lý các trường hợp hợp lệ.
\end{block}
\end{frame}

% 43
\begin{frame}[fragile]{Kiểm tra lựa chọn nhập vào}
\begin{lstlisting}[style=py]
choice = input("Continue? (Y/N): ")
choice = choice.upper()

if choice == "Y" :
    print("Continue")
elif choice == "N" :
    print("Stop")
else :
    print("Invalid choice")
\end{lstlisting}
\begin{exampleblock}{Mẹo}
Chuẩn hóa bằng \code{upper()} hoặc \code{lower()} để không phải kiểm tra cả chữ hoa và chữ thường.
\end{exampleblock}
\end{frame}

% 44
\begin{frame}{Hai đường tròn giao nhau: ý tưởng kiểm tra}
Để kiểm tra hai đường tròn có giao nhau, có thể so sánh khoảng cách giữa hai tâm với bán kính.
\begin{block}{Ba đại lượng chính}
\begin{itemize}
  \item khoảng cách giữa hai tâm: \code{distance};
  \item tổng bán kính: \code{r1 + r2};
  \item độ chênh bán kính: \code{abs(r1 - r2)}.
\end{itemize}
\end{block}
\begin{exampleblock}{Điều kiện giao nhau thông thường}
\code{abs(r1 - r2) <= distance and distance <= r1 + r2}
\end{exampleblock}
\end{frame}



% Added example + quiz
\begin{frame}[fragile]{Ví dụ bổ sung: kiểm tra số lượng mua hàng}
\begin{lstlisting}[style=py]
quantity = int(input("Quantity: "))

if quantity <= 0 :
    print("Quantity must be positive")
elif quantity > 100 :
    print("Quantity too large")
else :
    print("Accepted")
\end{lstlisting}
\begin{block}{Mẫu tư duy}
Kiểm tra lỗi trước, sau đó xử lý trường hợp hợp lệ.
\end{block}
\end{frame}

\begin{frame}{Quiz 3.7: input validation}
Một chương trình đọc số giờ làm việc trong tuần. Giá trị hợp lệ là từ 0 đến 80.
\begin{enumerate}
  \item Viết điều kiện phát hiện giá trị không hợp lệ.
  \item Nêu 5 ca kiểm thử nên dùng.
  \item Vì sao không nên tính lương trước khi kiểm tra số giờ hợp lệ?
\end{enumerate}
\end{frame}

\begin{frame}{Gợi ý đáp án Quiz 3.7}
\begin{enumerate}
  \item Điều kiện không hợp lệ: \code{hours < 0 or hours > 80}.
  \item Ca kiểm thử: \code{-1}, \code{0}, \code{1}, \code{80}, \code{81}.
  \item Vì dữ liệu sai có thể tạo kết quả sai hoặc gây lỗi ở các bước xử lý sau.
\end{enumerate}
\end{frame}

\section{Tổng kết và luyện tập}

% 45
\begin{frame}{Tổng kết: câu lệnh rẽ nhánh}
\begin{itemize}
  \item \code{if} chọn thực hiện một khối khi điều kiện đúng.
  \item \code{else} xử lý trường hợp điều kiện sai.
  \item \code{elif} phù hợp với nhiều phương án loại trừ lẫn nhau.
  \item Thụt lề xác định khối lệnh trong Python.
  \item Nên tránh lặp mã trong các nhánh.
\end{itemize}
\end{frame}

% 46
\begin{frame}{Tổng kết: điều kiện và Boolean}
\begin{itemize}
  \item Toán tử quan hệ tạo ra giá trị \code{True}/\code{False}.
  \item \code{==} kiểm tra bằng nhau; \code{=} là phép gán.
  \item Không nên so sánh bằng tuyệt đối với số thực sau tính toán.
  \item \code{and}, \code{or}, \code{not} kết hợp điều kiện.
  \item Đánh giá ngắn mạch có thể tránh lỗi và giảm tính toán.
\end{itemize}
\end{frame}

% 47
\begin{frame}{Tổng kết: kiểm thử và nhập liệu}
\begin{itemize}
  \item Mỗi nhánh quan trọng cần có ca kiểm thử.
  \item Giá trị biên thường là nơi dễ phát sinh lỗi.
  \item Kiểm tra dữ liệu đầu vào trước khi xử lý.
  \item Với văn bản, nên chuẩn hóa khi không cần phân biệt hoa/thường.
  \item Mô phỏng bằng tay giúp hiểu chương trình và tìm lỗi logic.
\end{itemize}
\end{frame}

% 48
\begin{frame}{Quiz nhanh}
\begin{enumerate}
  \item Sự khác nhau giữa \code{=} và \code{==} là gì?
  \item Khi nào nên dùng \code{elif} thay vì nhiều \code{if} độc lập?
  \item Vì sao điều kiện \code{score < 0 and score > 100} luôn sai?
  \item Giá trị biên của điều kiện \code{age >= 18} là những giá trị nào?
  \item \code{not (A and B)} tương đương với biểu thức nào?
\end{enumerate}
\end{frame}

% 49
\begin{frame}[fragile]{Bài tập trên lớp 1: phân loại điểm}
Viết chương trình đọc một điểm số nguyên và in ra:
\begin{itemize}
  \item \code{Invalid} nếu điểm ngoài khoảng 0--100;
  \item \code{A} nếu điểm từ 90 trở lên;
  \item \code{B} nếu điểm từ 80 đến dưới 90;
  \item \code{C} nếu điểm từ 70 đến dưới 80;
  \item \code{D} nếu điểm từ 60 đến dưới 70;
  \item \code{F} nếu điểm dưới 60.
\end{itemize}
\begin{block}{Yêu cầu}
Dùng \code{if/elif/else} và thiết kế ít nhất 6 ca kiểm thử.
\end{block}
\end{frame}

% 50
\begin{frame}[fragile]{Bài tập trên lớp 2: kiểm tra tên người dùng}
Viết chương trình đọc \code{username} và kiểm tra:
\begin{itemize}
  \item độ dài từ 6 đến 12 ký tự;
  \item không chứa khoảng trắng;
  \item chỉ gồm chữ cái và chữ số.
\end{itemize}
\smallgap
Nếu hợp lệ, in \code{Valid username}; ngược lại in thông báo lỗi phù hợp.
\smallgap
\begin{exampleblock}{Gợi ý}
Dùng \code{len()}, \code{" " in username}, \code{isalnum()}.
\end{exampleblock}
\end{frame}

% 51


% Added comprehensive quiz
\begin{frame}{Quiz tổng hợp bổ sung}
Trả lời nhanh.
\begin{enumerate}
  \item Điều kiện nào kiểm tra \code{x} nằm ngoài đoạn 1--10?
  \item Với \code{name = "An"}, vì sao \code{name == "an"} là \code{False}?
  \item Khi nào nên chuẩn hóa chuỗi bằng \code{upper()} hoặc \code{lower()}?
  \item Vì sao cần kiểm thử \code{49}, \code{50}, \code{51} với điều kiện \code{score >= 50}?
  \item \code{not (x < 0 or x > 100)} tương đương với điều kiện nào?
\end{enumerate}
\end{frame}

\begin{frame}{Đáp án quiz tổng hợp bổ sung}
\begin{enumerate}
  \item \code{x < 1 or x > 10}.
  \item Vì Python phân biệt chữ hoa và chữ thường.
  \item Khi không muốn phân biệt chữ hoa/chữ thường trong dữ liệu nhập.
  \item Vì đó là các giá trị ngay dưới, đúng tại, và ngay trên ngưỡng.
  \item \code{x >= 0 and x <= 100}.
\end{enumerate}
\end{frame}

\begin{frame}[fragile]{Ví dụ tổng hợp: phân loại đơn hàng}
\begin{lstlisting}[style=py]
member = input("Member? (Y/N): ").upper()
amount = float(input("Amount: "))

if amount < 0 :
    print("Invalid amount")
else :
    discount = 0
    if member == "Y" and amount >= 1000000 :
        discount = amount * 0.10
    elif member == "Y" :
        discount = amount * 0.05
    total = amount - discount
    print("Discount:", discount)
    print("Total:", total)
\end{lstlisting}
\end{frame}

\begin{frame}{Hướng dẫn học tập}
Sau bài này, khi gặp một bài toán có điều kiện, hãy luyện theo chu trình:
\begin{center}
\Large
\key{Hiểu bài toán} $\rightarrow$ \key{Xác định nhánh} $\rightarrow$ \key{Viết điều kiện} $\rightarrow$ \key{Kiểm thử biên}
\end{center}
\vspace{1em}
\begin{block}{Thói quen quan trọng}
Trước khi viết code, hãy tự hỏi: ``Có những trường hợp nào khác nhau? Mỗi trường hợp cần kiểm thử bằng dữ liệu nào?''
\end{block}
\end{frame}

\end{document}
